\documentclass{article}
\usepackage{../m}
\usepackage{graphicx}

\begin{document}
\noindent Paul Gustafson\\
\noindent Texas A\&M University - Math 467\\ 
\noindent Instructor: Stephen Fulling

\section{$\F_2$ and the Fano plane}
\subsection{Introduction}
The purpose of this paper is to answer Exercise 2.5 (p. 96) of Greenberg \cite{g}:
\begin{quote}
Let $\F_2$ be the field of two elements $\{0,1\}$, whose multiplication and addition have the usual tables except that $1+1 = 0$. Show that $\F_2^2$ is isomorphic to the smallest affine plane. Show that $P^2(\F_2)$ is isomorphic to the Fano plane.
\end{quote}

We will need a few preliminary definitions from Greenberg.

\begin{definition}
An \textbf{incidence geometry} $(\mathcal{P}, \mathcal{L}, \mathcal{I})$ consists of a set of points $\mathcal{P}$, a set of lines $\mathcal{L}$, and an incidence relation $\mathcal{I} \subset \mathcal{P} \cross \mathcal{L}$ such that:
\begin{enumerate}
\item Every pair of distinct points is incident to a unique line.
\item Every line is incident to at least two distinct points.
\item There exist three distinct noncollinear points. 
\end{enumerate}
\end{definition}

\begin{definition}
Two lines are $\textbf{parallel}$ if there is no point incident to both lines.
\end{definition}


\begin{definition}
A \textbf{projective plane} is an incidence geometry in which:
\begin{enumerate}
\item No two lines are parallel.
\item Every line is incident to at least three distinct points.
\end{enumerate}
\end{definition}

\begin{definition} 
An \textbf{affine plane} is an incidence geometry in which, for every line $l$ and point $P$ not incident to $l$, there exists a unique line $m$ incident to $P$ and parallel to $l$.
\end{definition}


\subsection{The affine plane $\F_2^2$}
As in $\R^2$, the points in $\F_2^2$ are simply the elements of the vector space $\F_2^2$, i.e. ordered pairs of elements of $\F_2$.  

Also analogously to $\R^2$, the lines in $\F_2^2$ are cosets of 1-dimensional subspaces of $\F_2^2$.  That is, every line in $\F_2^2$ can be written as $V + h$ for some 1-dimensional subspace $V\subset \F_2^2$ and $h \in \F_2^2$.  

Incidence in $\F_2^2$ corresponds to inclusion.  For example, the point $(1,1) \in \F_2^2$ is incident to the line $\{(1,0)t + (0,1): t \in \F_2\}$, since $(1,1) = (1,0) (1) + (0,1)$.

As Greenberg notes, the smallest affine plane, call it $\mathcal{A}$, consists of a set of four points $\{A,B,C,D\}$ and a set of four lines $\{\{A,B\}, \{B,C\}, \{C,D\}, \{D,A\}\}$, where incidence corresponds to inclusion. For example, the point $B$ is incident to the line $\{A,B\}$.

To see that $\mathcal{A}$ and $\F_2^2$ are isomorphic, first note that each 1-dimensional subspace over $\F_2$ has exactly 2 elements, so each line in $\F_2^2$ has 2 elements.  Conversely, given two elements $a,b \in \F_2^2$, the line $L((b-a)t, a)$ passes through $a$ and $b$.  Thus, the lines in $\F_2^2$ are precisely the two-element subsets of $\F_2^2$. 

Therefore, an arbitrary bijection $f$ from the points of $\F_2^2$ to the points of $\mathcal A$ induces a bijection of lines (two-element subsets), and since inclusion is preserved under bijections, incidence is also preserved. 

\subsection{$P^2(K)$}
For an arbitrary field $K$, the points of the projective space $P^2(K)$ are the 1-dimensional subspaces of $K^3$.  The lines are the 2-dimensional subspaces of $K^3$.  Incidence corresponds to containment. 

Projective points in $P^2(K)$ are usually denoted $(a\mathpunct{:} b \mathpunct{:} c)$ for some generator $(a,b,c) \in K^3\setminus \{0\}$. Then $(a\mathpunct{:}b\mathpunct{:}c) = (d\mathpunct{:}e\mathpunct{:}f)$ iff $(a,b,c)$ is a nonzero multiple of $(d,e,f)$. 

To label the projective lines, given an element $a \in K^3\setminus \{0\}$, consider the rank-1 linear transformation $T(a):K^3 \to K$ defined by $(T(a))(x) := \sum_{i=1}^3 a_ix_i$.  By the rank-nullity theorem, the nullity of $T(a)$ is 2. 

Moreover, if $V \subset K^3$ is a 2-dimensional subspace, then I claim $V = \ker(T(a))$ for some $a \in K^3\setminus \{0\}$.   Pick a basis $\{v, w\}$ for $V$. 
Since 
\begin{align*}
3 & \ge \dim(\ker(T(v)) + \ker(T(w))) 
\\ & = \dim(\ker(T(v)) + \dim(\ker(T(w)) - \dim(\ker(T(v)) \cap \ker(T(w)))  
\\ & = 4 - \dim(\ker(T(v)) \cap \ker(T(w))),
\end{align*}
 we have $\dim(\ker(T(v)) \cap \ker(T(w))) \ge 1$.

Hence, we may pick $a \in \ker(T(v)) \cap \ker(T(w)) \setminus \{0\}$.
Thus, if $x \in V$, then $x = \alpha v + \beta w$ for some $\alpha, \beta \in K$, so
\begin{align*}
(T(a))(x) & = \sum_{i=1}^3 a_i(\alpha v_i + \beta w_i)
\\ &  = \alpha \sum_{i=1}^3  v_i a_i + \beta \sum_{i=1}^3 w_i a_i 
\\ & = \alpha (T(v))(a) + \beta (T(w))(a) 
\\ & = 0.
\end{align*}
Therefore, $V \subset \ker(T(a))$, so by dimension counting, $V = \ker(T(a))$. Hence, the lines in $P^2(K)$, as 2-dimensional subspaces of $K^3$, are precisely the elements of the set $\{\ker(T(a)) : a \in K^3\setminus\{0\} \}$.

\begin{example}
The projective line $\{x+y+z=0: (x \mathpunct{:}y\mathpunct{:}z) \in P^2(\R)\}$ is incident to the point $(1\mathpunct{:}0\mathpunct{:}-1) \in P^2(\R)$ since $t+0-t = 0$ for all $t\in \R$.
\end{example}

\subsection{$P^2(\F_2)$ as the Fano plane}
A simplification occurs in $P^2(\F_2)$: there is a correspondence between each point in $P^2(\F_2)$ as a 1-dimensional subspace of $\F_2^3$ and its unique nonzero element in $\F_2^3$. Since, on the other hand, each non-zero element in $\F_2^3$ generates a 1-dimensional subspace of $F_2^3$, this correspondence defines a bijection from $P^2(\F_2)$ to $\F_2^3\setminus \{0\}$. Hence, there are $2^3 - 1 = 7$ points in $P^2(\F_2)$, and there is only one $(x \mathpunct{:}y\mathpunct{:}z)$-representation for each point.


Since every 2-dimensional subspace of $\F_2^3$ contains 0, a 1-dimensional subspace $V \subset \F_2^3$ lies within a 2-dimensional subspace $W \subset \F_2^3$ iff the unique nonzero element in $V$ lies within $W$. Hence, since each 2-dimensional subspace of $\F_2^3$ contains exactly $2^2 - 1 = 3$ nonzero $\F_2^3$-elements, each line in $P^2(\F_2)$ is incident to precisely 3 projective points.

Since we have seen that no two nonzero elements of an $\F_2$-vector space are linearly dependent, each pair of distinct nonzero elements in $\F_2^3$ determines a 2-dimensional subspace. Hence, since each 2-dimensional subspace of $\F_2^3$ contains exactly ${3 \choose 2} = 3$ distinct pairs of nonzero points, there are
$(1/3) {7 \choose 2} = 7$ dimension-2 subspaces of $\F_2^3$, i.e. projective lines in $P^2(\F_2)$.

From the previous section, each projective line can be written as $\ker(T(a))$ for some $a \in \F_2^3 \setminus\{0\}$.  Since there are 7 lines, the elements of $(\ker(T(a))_{a \in \F_2^3 \setminus\{0\}}$ must be distinct.  Thus, Figure \ref{fano} defines an explicit isomorphism between $P^2(\F_2)$ and the Fano plane.


\begin{figure}[h]
\centering
\includegraphics[width=3.5in]{Fano_plane.png}
\caption{An isomorphism between $P^2(\F_2)$ and the Fano plane}
\label{fano}
\end{figure}

%168 automorphisms of the Fano plane, what group?
However, this isomorphism is far from unique: the group of automorphisms of the Fano plane has order 168. Indeed, to see that that the order of this group is at most 168, first pick any 3 noncollinear points.  Under any automorphism, these 3 points must map to 3 noncollinear points, so there are $(7)(6)(4) = 168$ choices for the images of these three points. However, each pair of these points determines a distinct line, and the sole other point on that line must remain collinear with the pair.  Hence, since there are 3 such pairs, the images of 3 more points are fixed.  But there are only 7 points, so the last point's image is also determined.

Conversely, consider the action of $GL(3, \F_2)$, the group of nonsingular linear transformations of $\F_2^3$, on $P^2(\F_2)$. This group action is well-defined since the action of $GL(3, \F_2)$  on $F_2^3$ preserves subspace dimension and inclusion.  Moreover, if $g \in GL(3, \F_2)$ fixes $P^2(\F_2)$, then, since every 1-dimensional subspace of $\F_3$ has only one nonzero point, $g$ must fix $\F_2^3$. Thus, $GL(3, \F_2)$ acts faithfully on $P^2(\F_2)$, so is isomorphic to a subgroup of the automorphism group of the Fano plane.   Finally, by counting column choices, $|GL(3, \F_2)| = (7)(6)(3) = 168$.

\bibliographystyle{plain}
\bibliography{paper1}
\end{document}
